\documentclass[10pt]{article}
\usepackage{amsmath, amssymb, amsthm}
\usepackage{geometry}
\geometry{a4paper, margin=1in}

\usepackage{lmodern}
\usepackage{parskip}
\usepackage{enumitem}
\usepackage{hyperref}
\usepackage{multicol}
\usepackage{xcolor}
\usepackage{fancyhdr}
\usepackage{graphicx}
\usepackage{amsmath, amssymb, amsthm}
\usepackage{tikz}
\usetikzlibrary{calc}
\usepackage{xfrac}
\usetikzlibrary{calc}

\definecolor{mypink}{RGB}{255, 105, 180}
\definecolor{mydarkpink}{RGB}{199, 21, 133}
\definecolor{mygray}{RGB}{80, 80, 80}
\hypersetup{
    colorlinks=true,
    linkcolor=mydarkpink,
    urlcolor=mydarkpink,
    citecolor=mydarkpink
}


\def\HS{\hspace{\fontdimen2\font}}

\pagestyle{fancy}
\fancyhf{}
\fancyhead[L]{\textcolor{mygray}{Elijah's Math Notes}}
\fancyhead[R]{\textcolor{mygray}{\thepage}}
\fancyfoot[C]{\textcolor{mygray}{Prepared by Elijah}}

\usepackage{titling}
\pretitle{\begin{center}\Large\color{mypink}}
\posttitle{\par\end{center}\vskip 0.5em}
\preauthor{\begin{center}\large\color{mygray}}
\postauthor{\par\end{center}}
\predate{\begin{center}\small\color{mygray}}
\postdate{\par\end{center}}

\theoremstyle{plain}
\newtheorem{theorem}{Theorem}[section]
\newtheorem{lemma}[theorem]{Lemma}
\newtheorem{proposition}[theorem]{Proposition}
\newtheorem{corollary}[theorem]{Corollary}

\theoremstyle{definition}
\newtheorem{definition}[theorem]{Definition}
\newtheorem{example}[theorem]{Example}

\theoremstyle{remark}
\newtheorem{remark}[theorem]{Remark}
\newtheorem{note}[theorem]{Note}

\newcommand{\R}{\mathbb{R}}
\newcommand{\N}{\mathbb{N}}
\newcommand{\Z}{\mathbb{Z}}
\newcommand{\todo}[1]{\textcolor{mydarkpink}{\textbf{TODO:} #1}}
\newcommand{\highlight}[1]{\textcolor{mypink}{#1}}

\renewcommand{\theenumi}{\Alph{enumi}}

\definecolor{subtleRed}{RGB}{255, 66, 66}

\title{Vermont State Mathematics Coalition Talent Search Test \#3}
\author{Elijah Renner}
\date{\today}

\begin{document}

\maketitle
\tableofcontents

\section{Problem \#1} 

\paragraph{Problem} Note that in chess, a bishop attacks all squares that share a diagonal with it. An 8 x 8
chessboard has its squares labeled using the integers 1 through 64 inclusive as shown here.
\[
\begin{array}{|c|c|c|c|c|c|c|c|}
\hline
1 & 2 & 3 & 4 & 5 & 6 & 7 & 8 \\ \hline
9 & 10 & 11 & 12 & 13 & 14 & 15 & 16 \\ \hline
17 & 18 & 19 & 20 & 21 & 22 & 23 & 24 \\ \hline
25 & 26 & 27 & 28 & 29 & 30 & 31 & 32 \\ \hline
33 & 34 & 35 & 36 & 37 & 38 & 39 & 40 \\ \hline
41 & 42 & 43 & 44 & 45 & 46 & 47 & 48 \\ \hline
49 & 50 & 51 & 52 & 53 & 54 & 55 & 56 \\ \hline
57 & 58 & 59 & 60 & 61 & 62 & 63 & 64 \\ \hline
\end{array}
\]

\begin{enumerate}[label=(\alph*)]
	\item If \(n\) is divisible by 2,3,5, or 7, show that there exists a \(2n\)-digit integer \(N\) with nonzero digits such that no chop of \(N\) is divisible by \(n\).
	\item If \(N\) is a \(2n\)-digit integer with nonzero digits, and \(n\) is not divisible by 2,3,5, or 7, show that there is some chop of \(N\) that is divisible by \(n\).
\end{enumerate}

\paragraph{(a) Solution}

\paragraph{(b) Solution}

\section{Problem \#2} 

\subsection{a}

\paragraph{Problem} For an integer \(n\), let \(S(n)\) denote the sum of the digits of \(n\) and let \(D(n)\) denote the total number of positive divisors of \(n\). For \(n=(10+1)(10^2+1)(10^4+1)...(10^{2^{2024}}+1)-1\), compute the value of \(D(S(D(D(S(n)))))\). 

\paragraph{Solution} 

\subsection{b}

\paragraph{Problem} Let \(A\) be the answer to part (a). Suppose \(x_1\) and \(x_2\) are the two real solutions to the equation \(10-4^{A^x}=4^{2-A^x}\), with \(x_1>x_2\). Compute \(\left(\sqrt{A}\right)^{x_1-x_2}\).

\paragraph{Solution} 

\subsection{c}

\paragraph{Problem} Let \(B\) be the answer to part (b). Point \(P\) lies inside square \(VMTS\) in such a way that the areas of \(PVM\), \(PMT\), and \(PTS\) are \(B\), \(3B\), and \(9B\), respectively. Find the area of triangle \(PSV\).

\paragraph{Solution} 

\section{Problem \#3}

\paragraph{Problem} Find the number of ordered pairs of nonnegative integers \((a,b)\) such that \(9a+6\sqrt{ab}+b+2025=702\sqrt{a}+234\sqrt{b}\).

\paragraph{Solution} Let \(x,y=\sqrt{a}, \sqrt{b}\) so that \(9x^2+6xy+y^2+2025=702x+234y\implies (3x+y)^2+2025=234(3x+y).\) Manipulating we have a quadtratic in \(3x+y\):\[(3x+y)^2-234(3x+y)+2025=0\]
then we can apply the quadratic formula to find \(3x+y=\frac{234\pm\sqrt{234^2-4(1)(2025)}}{2}}=9,225\). Then we're left with two cases. Either \(3x+y=9\) or \(3x+y=225\). If \(3x+y=9\), only four nonnegative integers \(x\) satisfy \(3x+y=9\) where \(y\) is also a nonnegative integer. Similarly, if \(3x+y=225\), it's clear that for \(y\) ranging from \(0\) to \(225\), only \((\frac{225}{3}=75)+1\text{ (for the case where } x=0)=76\) nonnegative integers \(x\) satisfy \(3x+y=225\). This gives \(4+76=\boxed{80}\) pairs \((x,y)=(\sqrt{a},\sqrt{b})\), and since these \(x,y\) are nonnegative integers, it follows that the 80 corresponding \((a,b)=(x^2,y^2)\) are also nonnegative integer pairs satisfying \(9a+6\sqrt{ab}+b+2025=702\sqrt{a}+234\sqrt{b}\).

\section{Problem \#4}

\paragraph{Problem} A polynomial \(p(x)\) of degree at most 2025 has the property that each of the values \(p(0),p(1),...,p(2025)\) is either 0 or 1. Compute the maximum possible value of \(p(2026)\).

\paragraph{Solution} 

\section{Problem \#5} 

\paragraph{Problem} Triangle \(ABC\) has \(AB=13, AC=14, BC=15\). Squares \(ABDE,BCHI,CAFG\) are constructed outside triangle \(ABC\), forming a hexagon \(DEFGHI\). Squares \(DEKJ,EFML,FGON,GHQP,HISR,IDUT\) are constructed outside hexagon \(DEFGHI\), forming a dodecagon \(JKLMNOPQRSTU\). Find the area of \(JKLMNOPQRSTU\). 

\paragraph{Solution} I solve this problem, in essence, by repeated use of the law of cosines. First, separate \(JKLMNOPQRSTU\) into the polygons \(AEKJDBCGONF\), \(AEKLMNF\), \(BDJUTSI\), and \(CHRQPOG\).

\paragraph{\(AEKJDBCGONF.\)} Note that the area of \(ABC=\sqrt{s(s-AB)(s-BC)(s-CA)}\) where \(s=\frac{AB+BC+CA}{2}\) by Heron's formula. So, \(A_{ABC}=84\). Then, outside \(ABC\) we have squares with areas known by construction: \(A_{ABDE}=13^2=169\), \(A_{BCHI}=15^2=225\), \(A_{CAFG}=14^2=196\). We know that the squares \(DEKJ\), \(FGON\), and \(HISR\) have the same side lengths as the squares they were attached to. Therefore, \(A_{DEKJ}=13^2=169\), \(A_{FGON}=14^2=196\), \(A_{HISR}=15^2=225\). Now that we have all components of \(AEKJDBCGONF\), we can find that \(A_{AEKJDBCGONF}=A_{ABC}+A_{ABDE}+A_{BCHI}+A_{CAFG}+A_{DEKJ}+A_{FGON}+A_{HISR}=84+169+225+196+169+196+225=1264\).

\paragraph{\(AEKLMNF.\)} Next we find the area of \(AEKLMNF\). First, by law of cosines \(\angle BAC=\arccos\left(\frac{AC^2+AB^2-BC^2}{2(AC)(AB)}\right)\therefore \angle EAF=360^{\circ}-2(90^{\circ})-\angle BAC=180^{\circ}-\angle BAC\). By law of cosines we also find that \(FE=\sqrt{EA^2+FA^2-2(EA)(FA)\cos(\angle EAF)}\). With \(FE\) law of cosines tells us that \(\angle EFA=\arccos\left(\frac{FE^2+FA^2-EA^2}{2(FE)(FA)}\right)\therefore \angle MFN=360^{\circ}-180^{\circ}-90^{\circ}-\angle EFA=90^{\circ}-\angle EFA\). Because \(EFML\) is a square, \(LE=MF=FE\). Then law of cosines says \(MN=\sqrt{NF^2+MF^2-2(NF)(MF)\cos(\angle MFN)}\). Then, since we have \(\angle EFA\) and \(\angle EAF\), we know that \(\angle FEA=180^{\circ}-\angle EAF -\angle EFA\implies \angle LEK=360^{\circ}-180^{\circ}-90^{\circ}-\angle FEA\) because \(\angle LEF\) is inside a square. With \(\angle LEK\) we can, again by law of cosines, find \(LK=\sqrt{KE^2+LE^2-2(LE)(KE)\cos(\angle LEK)}\). We then have 
\[A_{MFN}=\sqrt{s(s-MF)(s-NF)(s-MN)} \text{ where }s=\frac{MF+NF+MN}{2}\]
\[A_{EAF}=\sqrt{s(s-EA)(s-FA)(s-EF)} \text{ where }s=\frac{EA+FA+EF}{2}\]
\[A_{LEK}=\sqrt{s(s-LE)(s-KE)(s-LK)} \text{ where }s=\frac{LE+KE+LK}{2}\]
by Heron's formula and 
\[A_{FELM}=FE^2\]
Evaluating all generalizations above with the initial lengths of \(ABC\) and lengths \(NF=14,FA=14, EA=13, KE=13\) and summing the defined areas gives \(A_{AEKLMNF}=841.5\).

\paragraph{\(BDJUTSI\)} Next, we apply an identical procedure as described for \(AEKLMNF\) to find \(BDJUTSI\). First, by law of cosines \(\angle CBA=\arccos\left(\frac{BA^2+BC^2-CA^2}{2(BA)(BC)}\right)\therefore \angle IBD=360^{\circ}-2(90^{\circ})-\angle CBA=180^{\circ}-\angle CBA\). By law of cosines we also find that \(DI=\sqrt{IB^2+DB^2-2(IB)(DB)\cos(\angle IBD)}\). With \(DI\) law of cosines tells us that \(\angle IDB=\arccos\left(\frac{DI^2+DB^2-IB^2}{2(DI)(DB)}\right)\therefore \angle UDJ=360^{\circ}-180^{\circ}-90^{\circ}-\angle IDB=90^{\circ}-\angle IDB\). Because \(IDUT\) is a square, \(TI=UD=DI\). Then law of cosines says \(UJ=\sqrt{JD^2+UD^2-2(JD)(UD)\cos(\angle UDJ)}\). Then, since we have \(\angle IDB\) and \(\angle IBD\), we know that \(\angle DIB=180^{\circ}-\angle IBD -\angle IDB\implies \angle TIS=360^{\circ}-180^{\circ}-90^{\circ}-\angle DIB\) because \(\angle TID\) is inside a square. With \(\angle TIS\) we can, again by law of cosines, find \(TS=\sqrt{SI^2+TI^2-2(SI)(TI)\cos(\angle TIS)}\). We then have 
\[A_{UDJ}=\sqrt{s(s-UD)(s-JD)(s-UJ)} \text{ where }s=\frac{UD+JD+UJ}{2}\]
\[A_{IBD}=\sqrt{s(s-IB)(s-DB)(s-ID)} \text{ where }s=\frac{IB+DB+ID}{2}\]
\[A_{TIS}=\sqrt{s(s-TI)(s-SI)(s-TS)} \text{ where }s=\frac{TI+SI+TS}{2}\]
by Heron's formula and 
\[A_{DITU}=DI^2\]
Evaluating all generalizations above with the initial given lengths of \(ABC\) and lengths \(JD=13,DB=13, IB=15, SI=15\) and summing the defined areas gives \(A_{BDJUTSI}=972\).

\paragraph{\(CHRQPOG\)} Just like above. First, by law of cosines \(\angle ACB=\arccos\left(\frac{CB^2+CA^2-AB^2}{2(CB)(CA)}\right)\therefore \angle GCH=360^{\circ}-2(90^{\circ})-\angle ACB=180^{\circ}-\angle ACB\). By law of cosines we also find that \(HG=\sqrt{GC^2+HC^2-2(GC)(HC)\cos(\angle GCH)}\). With \(HG\) law of cosines tells us that \(\angle GHC=\arccos\left(\frac{HG^2+HC^2-GC^2}{2(HG)(HC)}\right)\therefore \angle QHR=360^{\circ}-180^{\circ}-90^{\circ}-\angle GHC=90^{\circ}-\angle GHC\). Because \(GHQP\) is a square, \(PG=QH=HG\). Then law of cosines says \(QR=\sqrt{RH^2+QH^2-2(RH)(QH)\cos(\angle QHR)}\). Then, since we have \(\angle GHC\) and \(\angle GCH\), we know that \(\angle HGC=180^{\circ}-\angle GCH -\angle GHC\implies \angle PGO=360^{\circ}-180^{\circ}-90^{\circ}-\angle HGC\) because \(\angle PGH\) is inside a square. With \(\angle PGO\) we can, again by law of cosines, find \(PO=\sqrt{OG^2+PG^2-2(OG)(PG)\cos(\angle PGO)}\). We then have 
\[A_{QHR}=\sqrt{s(s-QH)(s-RH)(s-QR)} \text{ where }s=\frac{QH+RH+QR}{2}\]
\[A_{GCH}=\sqrt{s(s-GC)(s-HC)(s-GH)} \text{ where }s=\frac{GC+HC+GH}{2}\]
\[A_{PGO}=\sqrt{s(s-PG)(s-OG)(s-PO)} \text{ where }s=\frac{PG+OG+PO}{2}\]
by Heron's formula and 
\[A_{HGPQ}=HG^2\]
Evaluating all generalizations above with the initial given lengths of \(ABC\) and lengths \(RH=13,HC=13, GC=15, OG=15\) and summing the defined areas gives \(A_{CHRQPOG}=1093.5\).

Recall from our separations that \(A_{JKLMNOPQRSTU}=A_{AEKJDBCGONF}+A_{AEKLMNF}+A_{BDJUTSI}+A_{CHRQPOG}=1264+841.5+972+1093.5=\boxed{4,171}\).

\section{Problem \#6}

\paragraph{Problem} Let \(g(n)=2n-1-\lfloor\lfloor n\sqrt{2}\rfloor \sqrt{2}\rfloor\), where \(\lfloor x\rfloor\) denotes the greatest integer \(\leq x\). Prove that for all positive integers \(n\), \[g(n)^2+g(n+1)^2+g(n+2)^2+g(n+3)^2+g(n+4)^2\]

equals either 1 or 2.

\paragraph{Proof}

\end{document}




\end{document}